\section{Related Work}
\label{sec:carp:rel_work}

Data structures such as LSM-Trees~\cite{ONeil1996}, B+ Trees~\cite{Comer1979:UbiquitousBTree}, and
their variants~\cite{Zhong:XRPInKernelStorage, Raju2017} are employed by online databases to order
data.
These data structures heavily reorganize data internally and are more inefficient at writes than
bulk post-processing approaches~\cite{Qiao2022:ClosingTreeVs}.
Multiple in-situ analysis frameworks have been proposed including PreDatA \cite{Zheng2010:PreDatA},
GLEAN \cite{Vishw2011:SimulationtimeDataAnalysis, Vishw2011:GLEAN}, NESSIE \cite{Oldfi2012:Nessie},
DataSpaces \cite{Docan2010:DataSpaces}, and ADIOS~\cite{Gu2018:QueryingLargeScientifica}.
These systems are designed such that auxiliary nodes are used to perform analysis tasks.
ADIOS supports in-situ generation of range query indices using FastBit~\cite{Wu2005:FastBit} (the
same bitmap index used by FastQuery) using auxiliary nodes.
In-situ generation of bitmap indices would be faster than post-processing via FastQuery, but at the
cost of dedicated resources, and the space overhead and query performance limitations of bitmap
indices would still remain.

Usher et al.~\cite{Usher2021:AdaptiveSpatiallyAware} describes an in-situ indexing system that
aggregates application data into spatial locality-preserving data layouts.
Their approach requires provisioning dedicated aggregator nodes which collect spatially-partitioned
particle data from the application, add a bitmap index to the data, and write it to storage.
Particle codes already have spatial partitions because of how they decompose the problem.
This system only preserves these pre-existing partitions.
CARP is able to create partitioning along arbitrary dimensions using all-to-all shuffling.
Further, CARP does not require dedicated resources, introduces more powerful distribution
monitoring constructs, overlaps indexing and I/O for extremely high storage utilization, and
confers maximum benefits of an in-situ partitioning approach.
Since CARP does not require any dedicated resources, these two approaches can be composed together
for richer partitioning capabilities.

Slalom~\cite{Olma2017:Slalom} and VETI~\cite{Marou2023:VETI} describe raw-data indexing approaches.
Raw data indexing works on the query path to exploit pre-existing order in data and adaptively
reorder data in response to query patterns.
Scalable variants of these techniques can complement approximate indexes on the write-path to
further optimize query performance for exploratory queries, but they are not a substitute for
in-situ indexing on the write path.
WiscKey~\cite{Lu2017:WiscKey} introduces the notion of separating keys and values in LSM-Trees for
lower write amplification, similar to our discussion on clustered vs auxiliary indexes.
SuRF~\cite{Zhang2018:SuRF} can help reduce redundant reads for range queries with sparse keyspaces
but can not help speed up queries for data with no keyspace gaps.

% \section{Conclusion}

% CARP demonstrates that a sorted, clustered layout for range queries over a parallel application can
% be generated in a single streaming pass without any impact on application runtime.
% This accelerates subsequent analyses regardless of the workflow employed --- it can provide highly
% efficient queries by itself, or accelerate a query plan with auxiliary indices, or provide tighter
% partitions for columnar storage formats to exploit.
% Due to its low memory footprint, CARP can co-exist with other in-situ approaches running on the
% same system.
% End-to-end analysis stacks that produce an optimal combination of in-situ and post-hoc indexes for
% a given set of query requirements remain future work.
